% !PRINTED 156
\weberpage{156}

\begin{center}
\S.\,48.

\gesperrt{Resultanten.}
\end{center}

Eine andere Anwendung der Theorie der symmetrischen Functionen, die übrigens die Discriminantenbildung als speciellen Fall enthält, ist die Bildung der Resultanten.

Es seien

\weq{1}{
f(x) = a_0\,x^n + a_1\,x^{n-1} + a_2\,x^{n-2} + \cdots + a_n,
}

\weq{2}{
\varphi(x) = b_0\,x^m + b_1\,x^{m-1} + b_2\,x^{m-2} + \cdots + b_m
}

zwei ganze rationale Functionen vom $n^{\mathrm{ten}}$ und $m^{\mathrm{ten}}$ Grade und $\alpha_1, \alpha_2 \ldots \alpha_n$ seien die Wurzeln der ersten, so dass auch

\weq{3}{
f(x) = a_0\,(x - \alpha_1)\,(x - \alpha_2) \ldots (x - \alpha_n)
}

gesetzt werden kann.

Nun ist das Product

\[
\varphi(\alpha_1)\,\varphi(\alpha_2) \ldots \varphi(\alpha_n)
\]

eine symmetrische Function der Wurzeln von $f(x)$ und kann also ganz und rational durch

\[
\frac{a_1}{a_0}, \quad \frac{a_2}{a_0} \ldots \frac{a_n}{a_0}
\]

ausgedrückt werden. Die Ordnung dieser symmetrischen Function ist $(m, m \ldots m)$ (\S.\,45) und daher ist

\weq{4}{
R = a_0^m\,\varphi(\alpha_1)\,\varphi(\alpha_2) \ldots \varphi(\alpha_n)
}

eine ganze rationale und homogene Function $m^{\mathrm{ter}}$ Ordnung von $a_0, a_1 \ldots a_n$. Sie ist auch, wie ihr Ausdruck zeigt, eine ganze homogene Function $n^{\mathrm{ter}}$ Ordnung von $b_0, b_1 \ldots b_m$. Sie wird die \gesperrt{Resultante} der beiden Functionen $f(x)$ und $\varphi(x)$ genannt. Diese Bezeichnung rechtfertigt sich durch eine zweite Darstellung.

Bezeichnen wir die Wurzeln von $\varphi(x)$ mit $\beta_1, \beta_2 \ldots \beta_m$, setzen also

\weq{5}{
\varphi(x) = b_0\,(x - \beta_1)\,(x - \beta_2) \ldots (x - \beta_m),
}

so wird nach (4)

\weq{6}{
R = a_0^m\,b_0^n \overset{i,\,k}{\Pi}(\alpha_i - \beta_k),
}

worin das Productzeichen $\Pi$ sich auf alle Indices $i = 1, 2 \ldots n$, $k = 1, 2 \ldots m$ bezieht,
%CONT>>
